\documentclass[../Odyssey_Project.tex]{subfiles}

\begin{document}
\setcounter{section}{12}
\section{Wedderburn Theory}

\subsubsection*{Setting and Goal}
\begin{itemize}
    \item Maschke (Corollary~\ref{cor:12.8}) shows that if $\mathrm{char}(F) \nmid |G|$ then every finitely generated $FG$-module decomposes as a direct sum of simple modules.
    \item This section identifies the finite-dimensional $F$-algebras with this property and uses the resulting structure theorems to finish Kolchin's theorem on unipotent subgroups.
    \item Throughout, all algebras are finite-dimensional $F$-algebras (with unit unless stated), and all modules are finitely generated, equivalently finite-dimensional as $F$-vector spaces (Lemma~\ref{lem:12.5}).
    \item The classification is due to J.~H.~M.~Wedderburn.
\end{itemize}

\subsubsection*{Semisimplicity for Modules}
\begin{itemize}
    \item Three equivalent formulations of semisimplicity for an $A$-module: every submodule is a direct summand; the module is a direct sum of simple submodules; the module is a sum of simple submodules.
    \item Semisimplicity is preserved under submodules and quotient modules.
    \item An algebra $A$ is \emph{semisimple} if every non-zero $A$-module is semisimple; equivalently, $A$ is semisimple as an $A$-module over itself.
\end{itemize}

\begin{lemma}
\label{lem:13.1}
The following statements about an $A$-module $M$ are equivalent:
\begin{enumerate}
    \item Any submodule of $M$ is a direct summand of $M$.
    \item $M$ is semisimple.
    \item $M$ is a sum (but not, a priori, a direct sum) of simple submodules.
\end{enumerate}
\end{lemma}
\begin{proof}
    We first prove that (1) implies (2). If $M = 0$, then there is nothing to prove. Otherwise, choose a non-zero submodule $S \leq M$ of minimal $F$-dimension. Then $S$ is simple. By (1), we may write $M = S \oplus M_1$ for some submodule $M_1$. The same direct-summand condition passes to $M_1$: if $L \leq M_1$, then $L$ is a direct summand of $M$, say $M = L \oplus W$, and hence
    \[
        M_1 = L \oplus (M_1 \cap W).
    \]
    Since $\dim_F(M_1) < \dim_F(M)$, induction on $\dim_F(M)$ gives that $M_1$ is a direct sum of simple submodules. Hence $M$ is also a direct sum of simple submodules, so $M$ is semisimple.\\
    The implication (2) implies (3) is immediate from the definition of semisimplicity. It remains to prove that (3) implies (1). Assume that $M$ is a sum of simple submodules, and let $N \leq M$. Choose a submodule $V \leq M$ maximal among those satisfying $N \cap V = 0$; such a $V$ exists because $M$ is finite-dimensional. We claim that $N + V = M$. Suppose not. Since $M$ is a sum of simple submodules, some simple submodule $S \leq M$ is not contained in $N + V$. Then $S \cap (N + V)$ is a proper submodule of the simple module $S$, so $S \cap (N + V) = 0$. In particular, $S \cap V = 0$, and hence $V \subsetneq V + S$. If $n \in N \cap (V + S)$, write $n = v + s$ with $v \in V$ and $s \in S$. Then $s = n - v \in S \cap (N + V) = 0$, so $n = v \in N \cap V = 0$. Thus $N \cap (V + S) = 0$, contradicting the maximality of $V$. Therefore $N + V = M$, and since $N \cap V = 0$, we have $M = N \oplus V$. Thus every submodule of $M$ is a direct summand.
\end{proof}

\begin{lemma}
\label{lem:13.2}
Submodules and quotient modules of semisimple modules are semisimple.
\end{lemma}
\begin{proof}
    Let $M$ be a semisimple $A$-module. We first show that quotient modules of $M$ are semisimple. Let $N \leq M$, and let $\pi \colon M \to M/N$ be the natural epimorphism. Since $M$ is semisimple, write $M = S_1 + \cdots + S_t$ as a sum of simple submodules. Then
    \[
        M/N = \pi(M) = \pi(S_1) + \cdots + \pi(S_t).
    \]
    Each $\pi(S_i)$ is isomorphic to a quotient of $S_i$, so it is either zero or simple. Hence $M/N$ is a sum of simple submodules, and Lemma~\ref{lem:13.1} shows that $M/N$ is semisimple.\\
    Now let $L \leq M$. By Lemma~\ref{lem:13.1}, $L$ is a direct summand of $M$, so $M = L \oplus K$ for some submodule $K$. The projection $M \to L$ is an epimorphism, and hence $L$ is isomorphic to a quotient of $M$. Since quotient modules of $M$ are semisimple by the preceding paragraph, $L$ is semisimple. Thus both submodules and quotient modules of semisimple modules are semisimple.
\end{proof}

\begin{lemma}
\label{lem:13.3}
The algebra $A$ is semisimple iff the $A$-module $A$ is semisimple.
\end{lemma}
\begin{proof}
    Suppose first that the $A$-module $A$ is semisimple, and let $M$ be a non-zero $A$-module generated by $m_1,\ldots,m_r$. The direct sum $A^r$ is semisimple, since it is a finite direct sum of semisimple modules. Define
    \[
        \varphi \colon A^r \to M, \qquad \varphi(a_1,\ldots,a_r) = a_1m_1 + \cdots + a_rm_r.
    \]
    This is an $A$-module epimorphism, so $M$ is isomorphic to a quotient module of the semisimple module $A^r$. By Lemma~\ref{lem:13.2}, $M$ is semisimple. Therefore every non-zero $A$-module is semisimple, and hence $A$ is a semisimple algebra.\\
    Conversely, if $A$ is a semisimple algebra, then every non-zero $A$-module is semisimple by definition. In particular, the regular $A$-module $A$ is semisimple.
\end{proof}

\subsubsection*{Counting Simple Summands}
\begin{itemize}
    \item Decomposing $A \cong S_1 \oplus \dots \oplus S_r$ into simple summands captures all isomorphism types of simple $A$-modules.
    \item Multiplicities of simple summands in a semisimple module are determined uniquely (\nameref{thm:jordan-holder}).
\end{itemize}

\begin{proposition}
\label{prop:13.4}
Let $A$ be a semisimple algebra, and suppose that as $A$-modules we have $A \cong S_1 \oplus \dots \oplus S_r$, where the $S_i$ are simple submodules of $A$. Then any simple $A$-module is isomorphic with some $S_i$.
\end{proposition}
\begin{proof}
    Let $\theta \colon S_1 \oplus \cdots \oplus S_r \to A$ be an $A$-module isomorphism.
    Let $S$ be a simple $A$-module. Choose $0 \neq s \in S$, and define
    \[
        \varphi \colon A \to S, \qquad \varphi(a) = as.
    \]
    This is an $A$-module homomorphism, since $\varphi(ba) = (ba)s = b(as) = b\varphi(a)$. Since $S$ is simple and $\varphi(1) = s \neq 0$, the image of $\varphi$ is a non-zero submodule of $S$, so $\varphi$ is surjective. For each $i$, let $\iota_i \colon S_i \to S_1 \oplus \cdots \oplus S_r$ be the natural inclusion and set $\varphi_i = \varphi\theta\iota_i$. If every $\varphi_i$ were zero, then $\varphi\theta$ would vanish on $S_1 \oplus \cdots \oplus S_r$, contrary to the fact that $\varphi$ is non-zero and $\theta$ is surjective. Hence $\varphi_i$ is non-zero for some $i$. By \nameref{lem:schur}, this non-zero homomorphism between simple modules is an isomorphism. Therefore $S \cong S_i$ for some $i$.
\end{proof}

\begin{proposition}
\label{prop:13.5}
Suppose that $A$ is a semisimple algebra, and let $S_1, \dots, S_r$ be a collection of simple $A$-modules such that every simple $A$-module is isomorphic with exactly one $S_i$. Let $M$ be an $A$-module, and write $M \cong n_1 S_1 \oplus \dots \oplus n_r S_r$ for some non-negative integers $n_i$. Then the $n_i$ are uniquely determined.
\end{proposition}
\begin{proof}
    The point is that the Jordan--H\"older theorem does not only record the length of a composition series; it records the simple factors, up to reordering and isomorphism. In the module
    \[
        n_1 S_1 \oplus \cdots \oplus n_r S_r,
    \]
    we can build a composition series by removing one simple summand at a time. This gives a composition series in which $S_i$ appears exactly $n_i$ times as a composition factor. Now suppose that the same module could also be written as
    \[
        M \cong m_1S_1 \oplus \cdots \oplus m_rS_r.
    \]
    This second decomposition gives another composition series, this time with $S_i$ appearing exactly $m_i$ times. By \nameref{thm:jordan-holder}, the two composition series have the same composition factors up to order and isomorphism. Since each simple module is isomorphic with exactly one of the $S_i$, the number of factors isomorphic to $S_i$ must be the same in both series. Thus $n_i=m_i$ for every $i$, and the multiplicities are uniquely determined.
\end{proof}

\subsubsection*{Matrix Algebras over Division Algebras}
\begin{itemize}
    \item For a finite-dimensional $F$-algebra $D$, the matrix algebra $\mathcal{M}_n(D)$ has dimension $n^2 \dim_F D$.
    \item $E_{ij}(\alpha)$ denotes the matrix with sole non-zero entry $\alpha$ in position $(i,j)$.
    \item $D^n$ is the natural left $\mathcal{M}_n(D)$-module of column vectors.
    \item A \emph{division algebra} is an algebra whose non-zero elements form a group under multiplication; field extensions of $F$ are division algebras, but non-commutative examples also exist.
\end{itemize}

\begin{theorem}
\label{thm:13.6}
Let $D$ be a division algebra, and let $n \in \N$. Then any simple $\mathcal{M}_n(D)$-module is isomorphic with $D^n$, and $\mathcal{M}_n(D)$ is isomorphic as $\mathcal{M}_n(D)$-modules with the direct sum of $n$ copies of $D^n$. In particular, $\mathcal{M}_n(D)$ is a semisimple algebra.
\end{theorem}
\begin{proof}
    First we show that $D^n$ is simple as a left $\mathcal{M}_n(D)$-module. Let $0 \neq U \leq D^n$, and choose a non-zero vector $u \in U$. Some entry of $u$ is non-zero; say the $j$th entry is $x \in D^\times$. Multiplying $u$ by $E_{jj}(x^{-1})$ gives the $j$th standard basis vector $e_j$, so $e_j \in U$. For every $i$, multiplying $e_j$ by $E_{ij}(1)$ gives $e_i$, so every standard basis vector lies in $U$. Finally, multiplying $e_i$ by $E_{ii}(\alpha)$ gives the vector with $\alpha$ in the $i$th position and zero elsewhere. Since any vector in $D^n$ is the sum of its coordinate vectors, this shows that $U$ contains every vector in $D^n$. Hence $U = D^n$, so $D^n$ is simple.\\
    Now let $C_k$ be the submodule of $\mathcal{M}_n(D)$ consisting of matrices whose only possibly non-zero entries are in the $k$th column. Left multiplication by matrices preserves columns, so each $C_k$ is an $\mathcal{M}_n(D)$-submodule, and
    \[
        \mathcal{M}_n(D) = C_1 \oplus \cdots \oplus C_n
    \]
    as $\mathcal{M}_n(D)$-modules. The map $C_k \to D^n$ sending a matrix to its $k$th column is an $\mathcal{M}_n(D)$-module isomorphism: left multiplication by a matrix sends the $k$th column to the same result as the usual action on that column. Hence $\mathcal{M}_n(D)$ is a direct sum of $n$ copies of the simple module $D^n$. By Lemma~\ref{lem:13.3}, $\mathcal{M}_n(D)$ is a semisimple algebra. Finally, Proposition~\ref{prop:13.4} shows that every simple $\mathcal{M}_n(D)$-module is isomorphic with one of these simple summands, and hence is isomorphic with $D^n$.
\end{proof}

\subsubsection*{Simple Algebras}
\begin{itemize}
    \item An algebra is \emph{simple} if its only two-sided ideals are itself and the zero ideal.
    \item Simple algebras are semisimple, and matrix algebras over division algebras are simple.
\end{itemize}

\begin{lemma}
\label{lem:13.7}
Simple algebras are semisimple.
\end{lemma}
\begin{proof}
    Let $\Sigma$ be the sum of all simple submodules of the regular left $A$-module $A$. This sum is non-zero, because $A$ is finite-dimensional and therefore contains a non-zero submodule of minimal dimension; such a submodule is simple, since any proper non-zero submodule would have smaller dimension. We claim that $\Sigma$ is a two-sided ideal. It is already a left ideal because it is an $A$-submodule of $A$. Now let $S$ be a simple submodule of $A$, and let $a \in A$. The map
    \[
        S \to A, \qquad s \mapsto sa
    \]
    is an $A$-module homomorphism, since $(bs)a = b(sa)$ for all $b \in A$. The image of a simple module under a homomorphism is isomorphic to a quotient of that simple module, so $Sa$ is either zero or simple. In either case $Sa \subseteq \Sigma$. Since every element of $\Sigma$ is a finite sum of elements from simple submodules, multiplying such a sum on the right by $a$ still lands in $\Sigma$. Hence $\Sigma a \subseteq \Sigma$ for every $a \in A$. Thus $\Sigma$ is also a right ideal, and hence a two-sided ideal.\\
    Because $A$ is simple and $\Sigma \neq 0$, we must have $\Sigma = A$. Thus $A$, as a left $A$-module, is a sum of simple submodules. By Lemma~\ref{lem:13.1}, the $A$-module $A$ is semisimple, and then Lemma~\ref{lem:13.3} shows that $A$ is a semisimple algebra.
\end{proof}

\begin{theorem}
\label{thm:13.8}
Let $D$ be a division algebra, and let $n \in \N$. Then $\mathcal{M}_n(D)$ is a simple algebra.
\end{theorem}
\begin{proof}
    Let $0 \neq J$ be a two-sided ideal of $\mathcal{M}_n(D)$. We show that $J = \mathcal{M}_n(D)$. Choose a non-zero matrix $M \in J$, and let $x$ be a non-zero entry of $M$, say in position $(r,s)$. Since $D$ is a division algebra, $x^{-1}$ exists. Because $J$ is a two-sided ideal,
    \[
        E_{sr}(x^{-1}) M E_{ss}(1) \in J.
    \]
    This product is the matrix $E_{ss}(1)$: the right multiplication by $E_{ss}(1)$ keeps only column $s$, and the left multiplication by $E_{sr}(x^{-1})$ moves row $r$ into row $s$ while multiplying the surviving entry $x$ by $x^{-1}$. Hence $E_{ss}(1) \in J$.\\
    For any $i,j$, we have
    \[
        E_{ij}(1) = E_{is}(1)E_{ss}(1)E_{sj}(1),
    \]
    so $E_{ij}(1) \in J$. Since $J$ is a right ideal, $E_{ij}(\alpha) = E_{ij}(1)E_{jj}(\alpha)$ also lies in $J$ for every $\alpha \in D$. The matrices $E_{ij}(\alpha)$ span $\mathcal{M}_n(D)$ as an $F$-vector space, so $J = \mathcal{M}_n(D)$. Thus the only two-sided ideals are $0$ and the whole algebra, and $\mathcal{M}_n(D)$ is simple.
\end{proof}

\subsubsection*{Direct Sums of Algebras}
\begin{itemize}
    \item The \emph{external direct sum} $B = B_1 \oplus \dots \oplus B_r$ has the Cartesian product of underlying sets with componentwise operations; each $B_i$ embeds as an ideal.
    \item A $B$-module $M$ becomes a $B_i$-module via $b_i m = (0,\dots,b_i,\dots,0) m$, and is simple/semisimple as a $B_i$-module iff it is simple/semisimple as a $B$-module.
    \item An \emph{internal direct sum} of ideals $B_1, \dots, B_r$ (vector-space direct sum) is isomorphic to the external direct sum.
    \item Two-sided ideals of a finite direct sum decompose componentwise.
\end{itemize}

\begin{lemma}
\label{lem:13.9}
Let $B = B_1 \oplus \dots \oplus B_n$ be a direct sum of algebras. Then the two-sided ideals of $B$ are exactly the sets of the form $J_1 \oplus \dots \oplus J_n$, where $J_i$ is a two-sided ideal of $B_i$ for each $i$.
\end{lemma}
\begin{proof}
    Let $J$ be a two-sided ideal of $B$, and set $J_i = J \cap B_i$ for each $i$. Each $J_i$ is a two-sided ideal of $B_i$: if $x \in J_i$ and $c \in B_i$, then $cx$ and $xc$ lie in $J$ because $J$ is a two-sided ideal of $B$, and they lie in $B_i$ because multiplication is componentwise. Clearly $J_1 \oplus \cdots \oplus J_n \subseteq J$. Conversely, take $b \in J$, and write
    \[
        b = b_1 + \cdots + b_n
    \]
    with $b_i \in B_i$. Let $e_i$ be the element of $B$ whose $i$th component is the identity of $B_i$ and whose other components are zero. This element acts as a projection onto the $i$th component. Since $J$ is a right ideal, $b_i = be_i$ lies in $J$, and therefore $b_i \in J_i$. Hence $b \in J_1 \oplus \cdots \oplus J_n$, proving that $J = J_1 \oplus \cdots \oplus J_n$.\\
    The converse is direct: if each $J_i$ is a two-sided ideal of $B_i$, then $J_1 \oplus \cdots \oplus J_n$ is closed under left and right multiplication by elements of $B$, since multiplication in $B$ is componentwise. Thus these are exactly the two-sided ideals of $B$.
\end{proof}

\begin{theorem}
\label{thm:13.10}
Let $r \in \N$. For each $1 \leq i \leq r$, let $D_i$ be a division algebra over $F$, let $n_i \in \N$, and let $B_i = \mathcal{M}_{n_i}(D_i)$. Let $B$ be the external direct sum of the $B_i$. Then $B$ is a semisimple algebra having exactly $r$ isomorphism classes of simple modules and exactly $2^r$ two-sided ideals, namely every sum of the form $\bigoplus_{j \in J} B_j$, where $J$ is a subset of $\{1, \dots, r\}$.
\end{theorem}
\begin{proof}
    By Theorem~\ref{thm:13.6}, each $B_i$ is a direct sum
    \[
        B_i \cong G_{i1} \oplus \cdots \oplus G_{in_i},
    \]
    where the $G_{ij}$ are simple $B_i$-modules and are all isomorphic to the natural module $D_i^{n_i}$. We regard each $G_{ij}$ as a $B$-module by letting $B$ act through its $i$th component. Explicitly, $(b_1,\ldots,b_r)$ acts on $G_{ij}$ as $b_i$ acts. This does not introduce any new submodules, so each $G_{ij}$ is still simple as a $B$-module. Since $B = B_1 \oplus \cdots \oplus B_r$, we obtain a decomposition of $B$ as a direct sum of simple $B$-modules. Therefore $B$ is semisimple by Lemma~\ref{lem:13.3}.\\
    Proposition~\ref{prop:13.4} now implies that every simple $B$-module is isomorphic with some $G_{ij}$. If $i=k$, then $G_{ij} \cong G_{k\ell}$ by Theorem~\ref{thm:13.6}. If $i \neq k$, they cannot be isomorphic as $B$-modules. Indeed, let $e_i=(0,\ldots,1_i,\ldots,0)$. This element acts as the identity on $G_{ij}$ but as zero on $G_{k\ell}$. If $f \colon G_{ij}\to G_{k\ell}$ were a $B$-module isomorphism, then for $v\neq0$ we would have
    \[
        f(v)=f(e_iv)=e_if(v)=0,
    \]
    a contradiction. Thus there are exactly $r$ isomorphism classes of simple $B$-modules.\\
    Finally, by Theorem~\ref{thm:13.8}, each $B_i$ has only the two two-sided ideals $0$ and $B_i$. Lemma~\ref{lem:13.9} therefore shows that every two-sided ideal of $B$ is obtained by choosing, for each $i$, either $0$ or $B_i$. Hence the two-sided ideals are precisely the sums $\bigoplus_{j \in J}B_j$ with $J \subseteq \{1,\ldots,r\}$, and there are $2^r$ of them.
\end{proof}

\subsubsection*{Endomorphism and Opposite Algebras}
\begin{itemize}
    \item For an $A$-module $M$, $\mathrm{End}_A(M)$ is an $F$-algebra under composition; we call it the \emph{endomorphism algebra} of $M$.
    \item The \emph{opposite algebra} $B^{op}$ has the same underlying additive structure but reversed multiplication: $a \cdot b := ba$. Note $(B^{op})^{op} = B$, and the opposite of a division algebra is again a division algebra.
    \item These two notions are linked by $B^{op} \cong \mathrm{End}_B(B)$: every $B$-endomorphism of $B$ is right multiplication by some element.
    \item Endomorphism algebras decompose along $A$-module direct sums into block form, and over a single simple module they yield a matrix algebra.
\end{itemize}

\begin{lemma}
\label{lem:13.11}
Let $B$ be an algebra. Then $B^{op} \cong \mathrm{End}_B(B)$.
\end{lemma}
\begin{proof}
    For $a \in B$, let $P_a \colon B \to B$ be right multiplication by $a$, so $P_a(x) = xa$. This is a left $B$-module homomorphism because
    \[
        P_a(bx) = (bx)a = b(xa) = bP_a(x).
    \]
    Conversely, if $\varphi \in \mathrm{End}_B(B)$, then for every $b \in B$ we have
    \[
        \varphi(b) = \varphi(b \cdot 1) = b\varphi(1).
    \]
    Thus $\varphi = P_{\varphi(1)}$, so every $B$-endomorphism of the regular module is right multiplication by a unique element of $B$.\\
    It remains to check the multiplication. In $B^{op}$ the product of $a$ and $b$ is $a \cdot b = ba$. Hence, for $x \in B$,
    \[
        (P_aP_b)(x) = P_a(xb) = xba = P_{ba}(x) = P_{a \cdot b}(x).
    \]
    Therefore $a \mapsto P_a$ is an algebra isomorphism $B^{op} \cong \mathrm{End}_B(B)$.
\end{proof}

\begin{lemma}
\label{lem:13.12}
Let $S_1, \dots, S_r$ be the distinct simple $A$-modules; for each $i$, let $U_i$ be a direct sum of copies of $S_i$, and let $U = U_1 \oplus \dots \oplus U_r$. Then we have $\mathrm{End}_A(U) \cong \mathrm{End}_A(U_1) \oplus \dots \oplus \mathrm{End}_A(U_r)$.
\end{lemma}
\begin{proof}
    Let $\varphi \in \mathrm{End}_A(U)$. We first show that $\varphi$ preserves each isotypic summand $U_i$. The image $\varphi(U_i)$ is isomorphic to a quotient of $U_i$, so every composition factor of $\varphi(U_i)$ is isomorphic to $S_i$. If $\varphi(U_i)$ were not contained in $U_i$, then its image in $U/U_i$ would be non-zero. That image would have a composition factor isomorphic to $S_i$. But $U/U_i$ is a direct sum of copies of the $S_j$ with $j \neq i$, so none of its composition factors is isomorphic to $S_i$. This contradiction shows that $\varphi(U_i) \subseteq U_i$ for every $i$.\\
    Thus restriction gives endomorphisms $\varphi_i = \varphi|_{U_i} \in \mathrm{End}_A(U_i)$, and we obtain a homomorphism
    \[
        \Gamma \colon \mathrm{End}_A(U) \to \mathrm{End}_A(U_1) \oplus \cdots \oplus \mathrm{End}_A(U_r),
        \qquad
        \Gamma(\varphi) = (\varphi_1,\ldots,\varphi_r).
    \]
    This map is injective because an endomorphism of $U$ is determined by its restrictions to the direct summands $U_i$. It is also surjective: given $(\psi_1,\ldots,\psi_r)$ with $\psi_i \in \mathrm{End}_A(U_i)$, define $\psi \in \mathrm{End}_A(U)$ by writing $u = u_1 + \cdots + u_r$ and setting
    \[
        \psi(u) = \psi_1(u_1) + \cdots + \psi_r(u_r).
    \]
    This gives $\Gamma(\psi) = (\psi_1,\ldots,\psi_r)$. Since composition is componentwise under this decomposition, $\Gamma$ is an algebra isomorphism.
\end{proof}

\begin{lemma}
\label{lem:13.13}
If $S$ is a simple $A$-module, then for any $n \in \N$ we have $\mathrm{End}_A(nS) \cong \mathcal{M}_n(\mathrm{End}_A(S))$.
\end{lemma}
\begin{proof}
    View $nS$ as the module of column vectors with entries in $S$. Given a matrix $\Phi = (\varphi_{ij}) \in \mathcal{M}_n(\mathrm{End}_A(S))$, define an endomorphism $\Gamma(\Phi)$ of $nS$ by the usual matrix rule:
    \[
        \Gamma(\Phi)
        \begin{pmatrix}
            s_1\\
            \vdots\\
            s_n
        \end{pmatrix}
        =
        \begin{pmatrix}
            \varphi_{11}(s_1)+\cdots+\varphi_{1n}(s_n)\\
            \vdots\\
            \varphi_{n1}(s_1)+\cdots+\varphi_{nn}(s_n)
        \end{pmatrix}.
    \]
    This is $A$-linear because each $\varphi_{ij}$ is $A$-linear. The assignment $\Gamma$ respects addition directly. For multiplication, if $\Phi=(\varphi_{ij})$ and $\Psi=(\psi_{ij})$, then the $i$th component of $\Gamma(\Phi)\Gamma(\Psi)(s_1,\ldots,s_n)^t$ is
    \[
        \sum_k \varphi_{ik}\Bigl(\sum_j \psi_{kj}(s_j)\Bigr)
        =
        \sum_j\Bigl(\sum_k \varphi_{ik}\psi_{kj}\Bigr)(s_j),
    \]
    which is exactly the $i$th component coming from the matrix product $\Phi\Psi$. Thus $\Gamma$ is an algebra homomorphism.\\
    If $\Gamma(\Phi)=0$, apply it to a vector whose only non-zero entry is an arbitrary $s \in S$ in the $j$th position. The $i$th component of the result is $\varphi_{ij}(s)$, so every $\varphi_{ij}$ is zero. Thus $\Gamma$ is injective. Conversely, let $\psi \in \mathrm{End}_A(nS)$. Let $\iota_j \colon S \to nS$ insert an element into the $j$th coordinate, and let $\pi_i \colon nS \to S$ be projection onto the $i$th coordinate. Then
    \[
        \psi_{ij} = \pi_i \psi \iota_j
    \]
    lies in $\mathrm{End}_A(S)$, and the matrix $(\psi_{ij})$ maps under $\Gamma$ to $\psi$. Hence $\Gamma$ is surjective, and therefore $\mathrm{End}_A(nS) \cong \mathcal{M}_n(\mathrm{End}_A(S))$.
\end{proof}

\begin{lemma}
\label{lem:13.14}
Suppose that $F$ is algebraically closed, and let $S$ be a simple $A$-module. Then $\mathrm{End}_A(S) \cong F$.
\end{lemma}
\begin{proof}
    By \nameref{lem:schur}, $\mathrm{End}_A(S)$ is a division algebra: every non-zero endomorphism of the simple module $S$ is invertible. Let $\varphi \in \mathrm{End}_A(S)$. If $\varphi=0$, then it is scalar multiplication by $0$. Otherwise, since $S$ is a finite-dimensional vector space over the algebraically closed field $F$, the $F$-linear map $\varphi$ has an eigenvalue $\lambda \in F$. Thus $\varphi - \lambda I$ has non-zero kernel, so it is not invertible. But $\varphi - \lambda I$ is still an element of the division algebra $\mathrm{End}_A(S)$; the only non-invertible element of a division algebra is $0$. Therefore $\varphi = \lambda I$.\\
    Hence every $A$-endomorphism of $S$ is scalar multiplication by an element of $F$. The map $F \to \mathrm{End}_A(S)$ given by $\lambda \mapsto \lambda I$ is therefore an algebra isomorphism.
\end{proof}

\begin{lemma}
\label{lem:13.15}
Let $B$ be an algebra. Then $\mathcal{M}_n(B)^{op} \cong \mathcal{M}_n(B^{op})$ for any $n \in \N$.
\end{lemma}
\begin{proof}
    Define
    \[
        \Psi \colon \mathcal{M}_n(B)^{op} \to \mathcal{M}_n(B^{op})
    \]
    by $\Psi(X) = X^t$. This map is clearly bijective and $F$-linear. We check multiplication carefully because both the transpose and the opposite multiplication reverse order. Let $X=(x_{ij})$ and $Y=(y_{ij})$ be elements of $\mathcal{M}_n(B)^{op}$. In the domain, $X \cdot Y$ means $YX$ using the ordinary multiplication in $\mathcal{M}_n(B)$. In the codomain, entries multiply in $B^{op}$, so $u \cdot v = vu$. Therefore
    \[
        (\Psi(X)\Psi(Y))_{ij}
        = \sum_k x_{ki} \cdot y_{jk}
        = \sum_k y_{jk}x_{ki}
        = (YX)_{ji}
        = \Psi(X \cdot Y)_{ij}.
    \]
    Thus $\Psi(X \cdot Y)=\Psi(X)\Psi(Y)$, and $\Psi$ is an algebra isomorphism.
\end{proof}

\subsubsection*{Wedderburn's Structure Theorems}
\begin{itemize}
    \item Wedderburn's main structure theorem (Theorem~\ref{thm:13.16}): a finite-dimensional algebra is semisimple iff it is isomorphic to a direct sum of matrix algebras over division algebras.
    \item Equivalently, semisimple iff a direct sum of simple algebras, confirming the consistency of terminology.
    \item An algebra is simple iff it is isomorphic to a single matrix algebra over a division algebra (Theorem~\ref{thm:13.17}).
    \item Over an algebraically closed field, every semisimple algebra is a direct sum of full matrix algebras over $F$ itself (Theorem~\ref{thm:13.18}).
\end{itemize}

\begin{theorem}
\label{thm:13.16}
The algebra $A$ is semisimple iff it is isomorphic with a direct sum of matrix algebras over division algebras.
\end{theorem}
\begin{proof}
    Suppose first that $A$ is semisimple. As a left $A$-module, decompose $A$ into its simple isotypic parts:
    \[
        A = U_1 \oplus \cdots \oplus U_r,
    \]
    where $U_i \cong n_iS_i$, the $S_i$ are pairwise non-isomorphic simple $A$-modules, and each $n_i$ is positive. By Lemma~\ref{lem:13.11},
    \[
        A^{op} \cong \mathrm{End}_A(A).
    \]
    By Lemma~\ref{lem:13.12}, the endomorphism algebra of $A$ splits across the isotypic parts:
    \[
        \mathrm{End}_A(A)
        \cong \mathrm{End}_A(U_1) \oplus \cdots \oplus \mathrm{End}_A(U_r).
    \]
    Since $U_i \cong n_iS_i$, Lemma~\ref{lem:13.13} gives
    \[
        \mathrm{End}_A(U_i) \cong \mathcal{M}_{n_i}(\mathrm{End}_A(S_i)).
    \]
    Combining these isomorphisms,
    \[
        A^{op}
        \cong
        \mathcal{M}_{n_1}(\mathrm{End}_A(S_1))
        \oplus \cdots \oplus
        \mathcal{M}_{n_r}(\mathrm{End}_A(S_r)).
    \]
    Taking opposite algebras and using Lemma~\ref{lem:13.15}, we obtain
    \[
        A
        \cong
        \mathcal{M}_{n_1}(\mathrm{End}_A(S_1)^{op})
        \oplus \cdots \oplus
        \mathcal{M}_{n_r}(\mathrm{End}_A(S_r)^{op}).
    \]
    By \nameref{lem:schur}, each $\mathrm{End}_A(S_i)$ is a division algebra, and hence so is its opposite algebra. Thus $A$ is isomorphic with a direct sum of matrix algebras over division algebras.\\
    Conversely, if $A$ is isomorphic with a direct sum of matrix algebras over division algebras, then Theorem~\ref{thm:13.10} shows that $A$ is semisimple.
\end{proof}

\begin{theorem}
\label{thm:13.17}
The algebra $A$ is simple iff it is isomorphic with a matrix algebra over a division algebra.
\end{theorem}
\begin{proof}
    Suppose first that $A$ is simple. By Lemma~\ref{lem:13.7}, $A$ is semisimple, so Theorem~\ref{thm:13.16} writes $A$ as a direct sum of $r$ matrix algebras over division algebras. By Theorem~\ref{thm:13.10}, such a direct sum has exactly $2^r$ two-sided ideals. Since $A$ is simple, it has exactly two two-sided ideals, namely $0$ and $A$. Hence $2^r=2$, so $r=1$. Therefore $A$ is isomorphic with a single matrix algebra over a division algebra.\\
    Conversely, every matrix algebra over a division algebra is simple by Theorem~\ref{thm:13.8}.
\end{proof}

\begin{theorem}
\label{thm:13.18}
Suppose that the field $F$ is algebraically closed. Then any semisimple algebra is isomorphic with a direct sum of matrix algebras over $F$.
\end{theorem}
\begin{proof}
    In the proof of Theorem~\ref{thm:13.16}, the division algebras that occur are the opposite algebras of the endomorphism algebras $\mathrm{End}_A(S_i)$, where the $S_i$ are simple $A$-modules. If $F$ is algebraically closed, Lemma~\ref{lem:13.14} gives $\mathrm{End}_A(S_i) \cong F$ for each $i$. Since $F^{op}=F$, the decomposition in Theorem~\ref{thm:13.16} becomes a direct sum of matrix algebras over $F$.
\end{proof}

\subsubsection*{Nilpotent Ideals and the Radical}
\begin{itemize}
    \item In this part, $A$ is an algebra possibly without unit. An element $x \in A$ is \emph{nilpotent} if $x^n = 0$ for some $n \in \N$.
    \item An ideal $I$ is \emph{nilpotent} if $I^n = 0$ for some $n$, where $I^n$ is the ideal spanned by all products of $n$ elements; an algebra is nilpotent if some such $n$ exists for $I = A$.
    \item A nilpotent algebra cannot have a unit.
    \item The sum of two nilpotent ideals is nilpotent, and (using finite-dimensionality) a unique largest nilpotent ideal exists.
    \item Submodules with semisimple quotient are closed under finite intersection, so a smallest such submodule of $A$ exists.
    \item For an $A$-module $M$, $\mathrm{Ann}(M) = \{a \in A \mid aM = 0\}$ is the \emph{annihilator} of $M$, an ideal of $A$; the \emph{common annihilator} of a family $\{M_t\}_{t \in \mathcal{T}}$ is $\bigcap_{t \in \mathcal{T}} \mathrm{Ann}(M_t)$.
    \item Wedderburn (Theorem~\ref{thm:13.23}): the largest nilpotent ideal, the common annihilator of all simple $A$-modules, and the smallest submodule of $A$ with semisimple quotient all coincide. This common object is the \emph{radical} $\mathrm{rad}(A)$, also called the Jacobson radical $J(A)$.
\end{itemize}

\begin{lemma}
\label{lem:13.19}
Let $I$ and $J$ be nilpotent ideals of $A$. Then $I + J$ is also a nilpotent ideal of $A$.
\end{lemma}
\begin{proof}
    Choose $m,n \in \N$ such that $I^m=0$ and $J^n=0$. Consider a product of $m+n$ elements of $I+J$. Expanding each factor as a sum of one element from $I$ and one element from $J$, the product becomes a sum of terms
    \[
        z_1z_2\cdots z_{m+n},
    \]
    where each $z_k$ lies in either $I$ or $J$. In any such term, either at least $m$ of the factors lie in $I$, or at least $n$ of the factors lie in $J$. If at least $m$ factors lie in $I$, group the product around those $I$-factors. The pieces between them stay inside $A$, and because $I$ is a two-sided ideal, multiplying an element of $I$ on either side by elements of $A$ still leaves an element of $I$. Thus the whole term lies in $I^m=0$. The case with at least $n$ factors in $J$ is the same. Hence every such term is zero, and therefore $(I+J)^{m+n}=0$. Thus $I+J$ is nilpotent.
\end{proof}

\begin{corollary}
\label{cor:13.20}
$A$ has a largest nilpotent ideal.
\end{corollary}
\begin{proof}
    Let $K$ be the sum of all nilpotent ideals of $A$. Since $A$ is finite-dimensional, this sum is already the sum of finitely many nilpotent ideals: choose finitely many elements spanning $K$, and each one lies in a finite sum of nilpotent ideals. Thus
    \[
        K = I_1 + \cdots + I_t.
    \]
    By repeated application of Lemma~\ref{lem:13.19}, this finite sum is nilpotent. By construction, $K$ contains every nilpotent ideal of $A$. Hence $K$ is the largest nilpotent ideal of $A$.
\end{proof}

\begin{lemma}
\label{lem:13.21}
Let $M$ and $N$ be submodules of $A$ such that $A/M$ and $A/N$ are semisimple $A$-modules. Then $A/(M \cap N)$ is a semisimple $A$-module.
\end{lemma}
\begin{proof}
    Define
    \[
        \Phi \colon A/(M \cap N) \to A/M \oplus A/N
    \]
    by
    \[
        \Phi(a + M \cap N) = (a+M, a+N).
    \]
    This is a well-defined $A$-module homomorphism. Its kernel consists of those $a + M \cap N$ such that $a \in M$ and $a \in N$, which is exactly the zero coset. Hence $\Phi$ is injective. Thus $A/(M \cap N)$ is isomorphic to a submodule of $A/M \oplus A/N$. The direct sum $A/M \oplus A/N$ is semisimple, and Lemma~\ref{lem:13.2} says that submodules of semisimple modules are semisimple. Therefore $A/(M \cap N)$ is semisimple.
\end{proof}

\begin{corollary}
\label{cor:13.22}
$A$ has a smallest submodule having semisimple quotient.
\end{corollary}
\begin{proof}
    Intersect all submodules $M \leq A$ for which $A/M$ is semisimple. Since $A$ is finite-dimensional, this intersection is already the intersection of finitely many such submodules: repeatedly intersecting with a new submodule strictly lowers dimension only finitely many times. Say these finitely many submodules are $M_1,\ldots,M_t$. By Lemma~\ref{lem:13.21} and induction on $t$, the quotient
    \[
        A/(M_1 \cap \cdots \cap M_t)
    \]
    is semisimple. Therefore this finite intersection itself has semisimple quotient. It is contained in every submodule of $A$ with semisimple quotient by construction, so it is the smallest such submodule.
\end{proof}

\begin{theorem}
\label{thm:13.23}
Let $A$ be an algebra with unit. Then the following are equal:
\begin{itemize}
    \item The largest nilpotent ideal of $A$.
    \item The common annihilator of all simple $A$-modules.
    \item The smallest submodule of $A$ having semisimple quotient.
\end{itemize}
\end{theorem}
\begin{proof}
    Let $I$ be the largest nilpotent ideal of $A$, let $\Lambda$ be the common annihilator of all simple $A$-modules, and let $M$ be the smallest submodule of the regular module $A$ having semisimple quotient. We prove
    \[
        I = \Lambda = M.
    \]
    First let $J$ be any nilpotent ideal of $A$, and let $S$ be a simple $A$-module. The subspace $JS$ is a submodule of $S$, so either $JS=0$ or $JS=S$. If $JS=S$, then $J^kS=S$ for every $k \geq 1$ by induction. This contradicts nilpotence of $J$, since $J^n=0$ for some $n$. Therefore $JS=0$. Thus every nilpotent ideal annihilates every simple module, so $I \subseteq \Lambda$.\\
    Conversely, take a composition series of the regular $A$-module:
    \[
        A=A_0 \supset A_1 \supset \cdots \supset A_r=0.
    \]
    Each quotient $A_k/A_{k+1}$ is simple, so $\Lambda$ annihilates it. Equivalently, $\Lambda A_k \subseteq A_{k+1}$ for every $k$. It follows inductively that $\Lambda^kA \subseteq A_k$. In particular, $\Lambda^rA=0$. Since $A$ has a unit, every element $x\in\Lambda^r$ satisfies $x=x1$, so $\Lambda^r\subseteq\Lambda^rA=0$. Hence $\Lambda$ is nilpotent, and therefore $\Lambda \subseteq I$. Thus $I=\Lambda$.\\
    Next, since $\Lambda$ annihilates every simple module, it annihilates every semisimple module, because a semisimple module is a direct sum of simple modules. The quotient $A/M$ is semisimple, so $\Lambda(A/M)=0$. In particular, for $\lambda\in\Lambda$ we have $\lambda(1+M)=\lambda+M=0$, and hence $\lambda\in M$. Thus $\Lambda \subseteq M$.\\
    It remains to prove $M \subseteq \Lambda$. Suppose, for contradiction, that $\Lambda \subsetneq M$. Then some element of $M$ does not annihilate all simple modules, so there is a simple $A$-module $S$ with $MS \neq 0$. Since $S$ is simple, $MS=S$. Hence there is some $s\in S$ with $Ms\neq0$; the set $Ms$ is a submodule of $S$, so simplicity gives $Ms=S$. We may therefore choose $m \in M$ such that $ms=-s$. Then $(m+1)s=0$, so $m+1$ lies in the annihilator of $s$:
    \[
        \mathrm{Ann}(s)=\{a \in A \mid as=0\}.
    \]
    This annihilator is a proper submodule of the regular module $A$, because $1s=s \neq 0$. Since $A$ is finite-dimensional, $\mathrm{Ann}(s)$ is contained in some maximal submodule $N$ of $A$. The quotient $A/N$ is simple, hence semisimple. By the defining minimality of $M$, we have $M \subseteq N$. Thus $m \in N$, and also $m+1 \in N$. Hence $1=(m+1)-m \in N$, forcing $N=A$, a contradiction. Therefore $M \subseteq \Lambda$. We conclude that $I=\Lambda=M$.
\end{proof}

\begin{corollary}
\label{cor:13.24}
Let $A$ be an algebra with unit. Then $A$ is semisimple iff $A$ has no non-zero nilpotent ideals.
\end{corollary}
\begin{proof}
    By Lemma~\ref{lem:13.3}, $A$ is semisimple iff the regular $A$-module $A$ is semisimple. This is equivalent to saying that the smallest submodule of $A$ having semisimple quotient is $0$. By Theorem~\ref{thm:13.23}, that smallest submodule is the radical, which is also the largest nilpotent ideal of $A$. Hence $A$ is semisimple iff the largest nilpotent ideal is $0$, equivalently iff $A$ has no non-zero nilpotent ideals.
\end{proof}

\subsubsection*{Algebras Without Unit; Nilpotence}
\begin{itemize}
    \item Any algebra (possibly without unit) embeds as an ideal in an algebra with unit, with quotient $F$.
    \item Consequence: an algebra without unit having radical zero cannot exist---no such algebras occur.
    \item Wedderburn: an algebra generated as an $F$-vector space by nilpotent elements is itself nilpotent. This finishes \nameref{thm:kolchin}.
\end{itemize}

\begin{lemma}
\label{lem:13.25}
Let $A$ be an algebra, possibly without unit. Then there exists an algebra $B$ with unit and an ideal $I$ of $B$ such that, as algebras, we have $I \cong A$ and $B/I \cong F$.
\end{lemma}
\begin{proof}
    Define $B$ to be the vector space $F \oplus A$, with multiplication
    \[
        (\lambda,a)(\mu,b) = (\lambda\mu,\lambda b+\mu a+ab).
    \]
    This multiplication is the usual multiplication on $A$ plus the missing scalar identity terms; expanding both $((\lambda,a)(\mu,b))(\nu,c)$ and $(\lambda,a)((\mu,b)(\nu,c))$ gives the same expression by associativity and distributivity in $A$. Thus $B$ is an algebra, and the element $(1,0)$ is a unit for $B$. Let
    \[
        I = \{(0,a) \mid a \in A\}.
    \]
    Then $I$ is an ideal of $B$, since multiplying $(0,a)$ on either side by an element of $B$ again gives an element whose first coordinate is zero. The map $A \to I$ given by $a \mapsto (0,a)$ is an algebra isomorphism. Finally, the projection $B \to F$ onto the first coordinate is an algebra epimorphism with kernel $I$, so $B/I \cong F$.
\end{proof}

\begin{theorem}
\label{thm:13.26}
Let $A$ be a non-zero algebra, possibly without unit. If $A$ has no non-zero nilpotent ideals, then $A$ is an algebra with unit.
\end{theorem}
\begin{proof}
    Let $B$ and $I$ be as in Lemma~\ref{lem:13.25}, so $I \cong A$ and $B/I \cong F$. Let $J$ be a nilpotent ideal of $B$. Then $J \cap I$ is a nilpotent ideal of $I$, and hence corresponds to a nilpotent ideal of $A$. By hypothesis, $A$ has no non-zero nilpotent ideals, so $J \cap I=0$.\\
    The natural map $B \to B/I$ restricts to an injection on $J$, and its image is $(I+J)/I$. Equivalently,
    \[
        J \cong J/(J\cap I) \cong (I+J)/I.
    \]
    This image is a nilpotent ideal of $B/I$. But $B/I \cong F$, and a field has no non-zero nilpotent ideals; the whole field is not nilpotent because it has a unit. Hence $(I+J)/I=0$, so $J=0$. Therefore $B$ has no non-zero nilpotent ideals. Since $B$ has a unit, Corollary~\ref{cor:13.24} shows that $B$ is semisimple.\\
    By Theorem~\ref{thm:13.16}, $B$ is a direct sum of matrix algebras over division algebras. Theorem~\ref{thm:13.10} describes the ideals of such a direct sum: every ideal is a direct sum of some of the matrix-algebra summands. In particular, every non-zero ideal of $B$ has a unit, namely the sum of the identity elements in the selected summands. Since $I \cong A$ is a non-zero ideal of $B$, $I$ has a unit. Transporting this unit across the isomorphism $A \cong I$ gives a unit for $A$.
\end{proof}

\begin{theorem}
\label{thm:13.27}
Suppose that the algebra $A$ can be generated as an $F$-vector space by a set consisting of nilpotent elements. Then $A$ is nilpotent.
\end{theorem}
\begin{proof}
    We first reduce to the case where $F$ is algebraically closed. Let $\overline{F}$ be an algebraic closure of $F$, and form the scalar extension
    \[
        \overline{A}=\overline{F}\otimes_F A.
    \]
    with multiplication $(\alpha\otimes a)(\beta\otimes b)=\alpha\beta\otimes ab$. Choosing an $F$-basis of $A$ shows that $A$ embeds in $\overline{A}$ by $a \mapsto 1 \otimes a$, and that $\dim_{\overline{F}}\overline{A}=\dim_F A$. If $x \in A$ is nilpotent, then $1 \otimes x$ is nilpotent in $\overline{A}$. Thus, if a set of nilpotent elements spans $A$ over $F$, the corresponding elements $1\otimes x$ span $\overline{A}$ over $\overline{F}$ and are still nilpotent. If $\overline{A}$ is nilpotent, then its embedded subalgebra $A$ is nilpotent. Hence it suffices to prove the theorem after replacing $F$ by $\overline{F}$; from now on assume that $F$ is algebraically closed.\\
    We argue by induction on $\dim_F A$. The zero algebra is nilpotent, and if $\dim_F A=1$, then $A$ is spanned by a nilpotent element, so a sufficiently long product of elements of $A$ is zero. Now assume $\dim_F A>1$. If $A$ has a non-zero nilpotent ideal $I$, then $A/I$ is generated by the images of the original nilpotent generators, so by induction $A/I$ is nilpotent. Say $(A/I)^m=0$, so $A^m \subseteq I$. Since $I$ is nilpotent, say $I^n=0$, we get $A^{mn} \subseteq I^n=0$. Hence $A$ is nilpotent.\\
    It remains to rule out the case where $A$ has no non-zero nilpotent ideals. By Theorem~\ref{thm:13.26}, $A$ has a unit, and by Corollary~\ref{cor:13.24}, $A$ is semisimple. Since $F$ is algebraically closed, Theorem~\ref{thm:13.18} gives
    \[
        A \cong \mathcal{M}_{n_1}(F) \oplus \cdots \oplus \mathcal{M}_{n_r}(F).
    \]
    Projecting the nilpotent generating set of $A$ onto any summand shows that each summand $\mathcal{M}_{n_i}(F)$ is spanned by nilpotent matrices. But over an algebraically closed field, a nilpotent matrix has only the eigenvalue $0$, so its trace is zero. Thus all nilpotent matrices lie in the proper subspace $\ker(\operatorname{tr})$ of $\mathcal{M}_{n_i}(F)$, and they cannot span the whole matrix algebra. This contradiction shows that the no-non-zero-nilpotent-ideal case cannot occur. Therefore $A$ is nilpotent.
\end{proof}

\begin{proposition}
\label{prop:13.28}
Let $V_n(F)$ be the space of column vectors of length $n \in \N$ with entries from a field $F$, and let $H$ be a unipotent subgroup of $\mathrm{GL}(n,F)$. Then there exists some $0 \neq \mathbf{v} \in V_n(F)$ such that $x\mathbf{v} = \mathbf{v}$ for all $x \in H$.
\end{proposition}
\begin{proof}
    Let $I$ denote the identity matrix, and let $W$ be the $F$-subspace of $\mathcal{M}_n(F)$ spanned by
    \[
        \{x-I \mid x \in H\}.
    \]
    If $x \in H$, then $x$ is unipotent, so $x-I$ is nilpotent. Also, for $x,y \in H$,
    \[
        (x-I)(y-I) = (xy-I) - (x-I) - (y-I).
    \]
    Since $xy \in H$, the right-hand side lies in $W$. Hence $W$ is closed under multiplication, so $W$ is an algebra, possibly without unit, generated as an $F$-vector space by nilpotent elements. By Theorem~\ref{thm:13.27}, $W$ is nilpotent.\\
    If $W=0$, then $x-I=0$ for every $x \in H$, so every non-zero vector is fixed by $H$. Otherwise, choose $t \in \N$ minimal such that every product
    \[
        (x_1-I)\cdots(x_t-I)
    \]
    with $x_1,\ldots,x_t \in H$ is zero. Such a $t$ exists because $W$ is nilpotent and the elements $x-I$ span $W$. By minimality, there exist $y_1,\ldots,y_{t-1} \in H$ such that
    \[
        P=(y_1-I)\cdots(y_{t-1}-I)
    \]
    is not the zero matrix. Choose a vector $\mathbf{v}'$ such that $\mathbf{v}=P\mathbf{v}' \neq 0$. For any $x \in H$, the product
    \[
        (x-I)(y_1-I)\cdots(y_{t-1}-I)
    \]
    has $t$ factors from $W$, and hence is zero. Therefore
    \[
        (x-I)\mathbf{v}=0
    \]
    for every $x \in H$, which is exactly $x\mathbf{v}=\mathbf{v}$ for every $x \in H$.
\end{proof}

\subsection*{Exercises}

Throughout these exercises, $A$ denotes a finite-dimensional algebra with unit over a field $F$, and all $A$-modules are finitely generated.

\begin{exercise}
\label{ex:13.1}
Let $n \in \N$. Let $V$ be the $n$-dimensional subspace of $A^n$ spanned by the identity elements of the summands. Show that if $M$ is an $A$-module and $T \colon V \to M$ is a linear transformation, then there is a unique $A$-module homomorphism from $A^n$ to $M$ which extends $T$.
\end{exercise}
\begin{solution}
\end{solution}

\begin{exercise}
\label{ex:13.2}
(cont.) Suppose that $B$ is an $A$-module having an $n$-dimensional subspace $U$ such that whenever $M$ is an $A$-module and $T \colon U \to M$ is a linear transformation, there is a unique extension of $T$ to an $A$-module homomorphism from $B$ to $M$. Show that $B$ and $A^n$ are isomorphic $A$-modules. (Modules of the form $A^n$ for some $n \in \N$ are called \emph{free} $A$-modules; hence this exercise provides an alternate characterization of freeness, which we shall generalize for group algebras in the exercises to Section 16.)
\end{exercise}
\begin{solution}
\end{solution}

\begin{exercise}
\label{ex:13.3}
Let $n \in \N$, and let $T_n(F)$ be the algebra of upper triangular $n \times n$ matrices. Show that the set $V_n(F)$ of column vectors over $F$ of length $n$ is a $T_n(F)$-module that has a unique composition series in which every simple $T_n(F)$-module appears exactly once as a composition factor.
\end{exercise}
\begin{solution}
\end{solution}

\begin{exercise}
\label{ex:13.4}
(cont.) Show that the $T_n(F)$-module $T_n(F)$ is isomorphic with the direct sum of all non-zero submodules of $V_n(F)$.
\end{exercise}
\begin{solution}
\end{solution}

\begin{exercise}
\label{ex:13.5}
(cont.) Find the largest nilpotent ideal of $T_n(F)$, and show without reference to Theorem~\ref{thm:13.23} that it is also the common annihilator of all simple $T_n(F)$-modules and the smallest submodule of $T_n(F)$ having semisimple quotient.
\end{exercise}
\begin{solution}
\end{solution}

\begin{exercise}
\label{ex:13.6}
Let $U$ be an $A$-module, let $n \in \N$, and let $U^n$ be the set of column vectors of length $n$ with entries from $U$, considered in the obvious way as an $\mathcal{M}_n(A)$-module. Show that $U$ is a simple $A$-module iff $U^n$ is a simple $\mathcal{M}_n(A)$-module.
\end{exercise}
\begin{solution}
\end{solution}

\begin{exercise}
\label{ex:13.7}
(cont.) Show that $\mathrm{Hom}_A(U,V) \cong \mathrm{Hom}_{\mathcal{M}_n(A)}(U^n, V^n)$ for any $A$-modules $U$ and $V$.
\end{exercise}
\begin{solution}
\end{solution}

\begin{exercise}
\label{ex:13.8}
(cont.) Show that every $\mathcal{M}_n(A)$-module is isomorphic with a module of the form $U^n$ for some $A$-module $U$. (What Exercises~\ref{ex:13.6}--\ref{ex:13.8} show is that the module theory of the algebras $A$ and $\mathcal{M}_n(A)$ is, in a sense that can be made precise, the same. In the language of category theory, we say that the categories of $A$-modules and $\mathcal{M}_n(A)$-modules are \emph{Morita equivalent}.)
\end{exercise}
\begin{solution}
\end{solution}

\begin{exercise}
\label{ex:13.9}
Show that $A$ is a simple $A$-module iff $A$ is a division algebra.
\end{exercise}
\begin{solution}
\end{solution}

\begin{exercise}
\label{ex:13.10}
We can clearly regard an $A/\mathrm{rad}\, A$-module as being an $A$-module; on the other hand, a simple $A$-module is annihilated by $\mathrm{rad}(A)$ and hence can be considered as an $A/\mathrm{rad}\, A$-module. Show that an $A$-module is simple iff it is a simple $A/\mathrm{rad}(A)$-module. (Since $A/\mathrm{rad}(A)$ is a semisimple algebra, this implies that the determination of simple modules over arbitrary algebras reduces to the case of semisimple algebras.)
\end{exercise}
\begin{solution}
\end{solution}

\begin{exercise}
\label{ex:13.11}
Let $B$ be a finite-dimensional $F$-algebra, possibly without unit. Show that if $B$ is a simple algebra whose multiplication is not identically zero, then $B$ is isomorphic with a matrix algebra over a division algebra.
\end{exercise}
\begin{solution}
\end{solution}

\subsection*{Further Exercises}

Let $A$ be as above, and let $M$ be a finitely generated $A$-module.

\begin{exercise}
\label{ex:13.12}
Show that the following objects (exist and) are equal: the intersection of all maximal submodules of $M$; the smallest submodule of $M$ having semisimple quotient; and $\mathrm{rad}(A)M$. This submodule of $M$ is called the \emph{radical} of $M$ and is denoted $\mathrm{rad}(M)$. (Observe that our two definitions of the symbol $\mathrm{rad}(A)$ agree.)

We define a descending series $\mathrm{rad}^n(M)$ of submodules of $M$ by setting $\mathrm{rad}^0(M) = M$ and $\mathrm{rad}^n(M) = \mathrm{rad}(\mathrm{rad}^{n-1}(M))$ for $n \in \N$. It follows from Exercise~\ref{ex:13.12} that $\mathrm{rad}^n(M) = (\mathrm{rad}(A))^n M$ for any $n \in \N$. This series is called the \emph{radical series} of $M$. Observe that $\mathrm{rad}^n(M) = 0$ for some $n$, since $\mathrm{rad}(A)$ is a nilpotent ideal, and that each successive quotient of the radical series is a semisimple module, with $M$ being semisimple iff $\mathrm{rad}(M) = 0$.
\end{exercise}
\begin{solution}
\end{solution}

\begin{exercise}
\label{ex:13.13}
(cont.) Show that $M$ has a unique largest semisimple submodule and that this submodule is equal to $\{m \in M \mid \mathrm{rad}(A)m = 0\}$. We call this submodule the \emph{socle} of $M$, and we denote it by $\mathrm{soc}(M)$.
\end{exercise}
\begin{solution}
\end{solution}

\begin{exercise}
\label{ex:13.14}
(cont.) We define an ascending series $\mathrm{soc}^n(M)$ of submodules of $M$ as follows: We let $\mathrm{soc}^0(M) = 0$ and $\mathrm{soc}^1(M) = \mathrm{soc}(M)$, and for $n > 1$ we let $\mathrm{soc}^n(M)$ be the submodule of $M$ such that $\mathrm{soc}^n(M)/\mathrm{soc}^{n-1}(M) = \mathrm{soc}(M/\mathrm{soc}^{n-1}(M))$. Show that we have $\mathrm{soc}^n(M) = \{m \in M \mid \mathrm{rad}(A)^n m = 0\}$ for any $n \in \N$. (This series is called the \emph{socle series} of $M$. As with the radical series, we observe that $\mathrm{soc}^n(M) = M$ for some $n$, since $\mathrm{rad}(A)$ is nilpotent, and that each successive quotient of the socle series is semisimple, with $M$ being semisimple iff $\mathrm{soc}(M) = M$.)
\end{exercise}
\begin{solution}
\end{solution}

\begin{exercise}
\label{ex:13.15}
(cont.) Show that $\mathrm{rad}^n(M) = 0$ iff $\mathrm{soc}^n(M) = M$; conclude that the radical and socle series of $M$ have the same finite length. If this common length is $r$, then show that $\mathrm{rad}^n(M) \subseteq \mathrm{soc}^{r-n}(M)$ for every $0 \leq n \leq r$.
\end{exercise}
\begin{solution}
\end{solution}

\end{document}
